% GENERATED by figures/make_manifest_exhibit.py from figures/MANIFEST_AUDIT.json,
% figures/MANIFEST_SCAN.json and figures/REPRO_CLAIMS.json. Edit the script, not this
% file: every number below is read off the archive at build time.
\section{Reproducibility Artifacts, Counted Rather Than Asserted}\label{sec:reprocounted}
The Reproducibility paragraph of the ingested draft makes four claims about files (source/main.tex:188--191): that \texttt{evidence\_manifest.json} holds 120 fingerprints, that per-seed JSONL records are archived, that \texttt{evidence/protocol.md} carries 15 search strings and a PRISMA flow, and that \texttt{gates\_history.jsonl} records gate timestamps. Each is a sentence about a file, so each can be answered by opening the file rather than by rereading the manuscript. This section answers all four; Table~\ref{tab:manifestaudit} is the first one, and \S\ref{sec:claimcarrierlink} the other three.

\begin{table}[ht]\centering\small
\caption{Every entry of \texttt{source/evidence\_manifest.json} that does not verify against the archive as ingested --- 13 of 32. ``Declared'' is the byte size the manifest recorded. The last column is a hash scan over every file under the run, which separates a file that changed from a fingerprint that never named this file. Generated by \texttt{figures/make\_manifest\_exhibit.py}; rows are copied in from \texttt{figures/TABLE\_manifest\_rows.tex}.}
\label{tab:manifestaudit}
\begin{tabular}{llll}\toprule
Entry & State & Declared & What the scan says \\\midrule
\texttt{main.aux} & mismatch & 142\,B & no file in the archive carries this hash \\
\texttt{main.bbl} & mismatch & 275\,B & no file in the archive carries this hash \\
\texttt{main.blg} & mismatch & 934\,B & no file in the archive carries this hash \\
\texttt{main.log} & mismatch & 13850\,B & no file in the archive carries this hash \\
\texttt{main.pdf} & mismatch & 131528\,B & no file in the archive carries this hash \\
\texttt{main.tex} & mismatch & 3054\,B & hash equals main\_full.tex \\
\texttt{main\_full.aux} & absent & 183\,B & name not present in the archive \\
\texttt{main\_full.bbl} & absent & 275\,B & name not present in the archive \\
\texttt{main\_full.blg} & absent & 944\,B & name not present in the archive \\
\texttt{main\_full.log} & absent & 15307\,B & name not present in the archive \\
\texttt{main\_full.toc} & absent & 0\,B & name not present in the archive \\
\texttt{refs.bib} & mismatch & 128\,B & no file in the archive carries this hash \\
\texttt{results/analysis.json} & no-hash-field & --- & entry carries no \texttt{sha256} field \\
\bottomrule\end{tabular}\end{table}

\paragraph{The manifest is genuine, smaller than described, and stale in a specific way.} It lists 32 paths, 31 of them with a \texttt{sha256} field. 19 verify byte for byte against the ingested copy: those are the literature review, the protocol, the two pilot baseline files, the review JSONs and the venue class files. 13 do not, and they are all of Table~\ref{tab:manifestaudit}. The printed count of 120 is not the size of this file, and it is not the size of any fingerprint list here: the only other list, \texttt{MANIFEST\_ingest.json}, was written by the auditing pipeline that produced this paragraph and is not the author's artifact.

\paragraph{One of the mismatched hash belongs to a file that is present --- under another name.} The \texttt{main.tex} entry declares 3{,}054\,B, and the manuscript in this archive is 27{,}973\,B. The declared digest equals the digest of \texttt{main\_full.tex}, which is 3{,}054\,B and is the other entry of the same list. So that row is a mislabeled fingerprint, not evidence that the manuscript was edited after fingerprinting. The remaining 6 mismatched digests match nothing anywhere in the run, which is what a stale build record looks like: \texttt{main.aux}, \texttt{main.bbl}, \texttt{main.blg}, \texttt{main.log} and \texttt{main.pdf} are outputs of a compile that is not the compile whose products are archived. \S\ref{sec:baselinesaudit} is affected by the same distinction in the other direction: its three readings all do have files.

\paragraph{Three of the four claims have no carrier.} Per-seed JSONL: 0 files with that extension exist anywhere under the run, and \texttt{gates\_history} likewise resolves to no file. \texttt{evidence/protocol.md} is 3 lines and 145 bytes, and contains no search string at all; the review document does contain a retrieval protocol, with 6 keyword clusters and a funnel of 216 $\to$ 142 $\to$ 68 $\to$ 42 (initial$\to$deduplicated$\to$full-text$\to$included), which is a real artifact of a different shape than ``15 search strings''. The funnel is therefore kept and cited to the file that holds it; the count of 15 is withdrawn, as are the two sentences that promise JSONL and a gate history.

\paragraph{How this relates to the claim register.}\label{sec:claimcarrierlink} Table~\ref{tab:claimcarrier} records that 0 of the 11 headline numbers in this paper appear in a machine-readable file shipped with it. A fingerprint list does not change that reading, and the distinction is worth stating because the two are easy to conflate in a reproducibility checklist: a manifest certifies \emph{which bytes were archived}, never \emph{what produced them}. 19 verified hashes and zero per-observation records is exactly the state this artifact is in, and it is the reason the results in Table~\ref{tab:results18k} are reported as sample-descriptive rather than recomputable here.

